% =============================================================================
% CAPÍTULO 5 - IMPLEMENTACIÓN
% =============================================================================
% Enlaza con el capítulo de Diseño (cap.~\ref{chap:diseno}) sin repetir su contenido.
% =============================================================================

\chapter{Implementación}
\label{chap:implementacion}

Este capítulo describe cómo se ha llevado a la práctica el diseño definido en el Capítulo~\ref{chap:diseno}: las decisiones de implementación más relevantes en el frontend y el despliegue en producción, con especial atención a la persistencia de datos y del modelo en el servidor. Se explica el \textit{cómo} y el \textit{por qué} de cada elección de forma breve, sin volver a detallar la arquitectura ni el modelo de datos ya presentados.

% -----------------------------------------------------------------------------
\section{Implementación del frontend}
% -----------------------------------------------------------------------------

La interfaz de usuario se ha implementado siguiendo la estructura en capas y rutas descrita en la sección~\ref{sec:estructura-frontend}. A continuación se resumen las decisiones que más impacto tienen en el comportamiento y el mantenimiento de la aplicación.

\subsection{Capa de servicios y comunicación con el backend}

Se centralizaron todas las llamadas HTTP al backend en un único servicio (\texttt{api.ts}) que expone métodos por recurso (mazos, partidas, cartas, favoritos, recomendaciones, etc.). \textbf{Por qué:} (1) una sola definición de la URL base del API, que en producción se configura mediante la variable de entorno \texttt{VITE\_API\_BASE\_URL} para apuntar al servicio desplegado en Render; (2) envío consistente de la cabecera \texttt{X-Auth0-ID} cuando el usuario está autenticado, necesaria para que el backend identifique al usuario; (3) reutilización de la lógica de manejo de errores y de respuestas JSON. La comunicación con MarvelCDB se delegó en \texttt{marvelcdbService.ts} para mantener separada la lógica de importación externa.

\subsection{Autenticación y rutas protegidas}

La autenticación se implementó con Auth0 mediante el SDK para React (\texttt{@auth0/auth0-react}). El estado de sesión (usuario, login, logout) se expone en un \texttt{AuthContext} que envuelve la aplicación. \textbf{Por qué:} así cualquier componente puede saber si hay usuario logueado y mostrar u ocultar opciones (p. ej. ``Mis mazos'', ``Favoritos'') sin duplicar lógica. Las rutas que requieren sesión (crear/editar mazo, perfil, recomendación IA, etc.) se envuelven en el componente \texttt{ProtectedRoute}, que comprueba la autenticación antes de renderizar y redirige al login de Auth0 si no hay sesión. \textbf{Por qué:} se evita mostrar contenido privado por error y se centraliza la protección en un único punto.

\subsection{Layout común y navegación SPA}

Todas las páginas comparten un \texttt{Layout} que incluye cabecera y pie. La cabecera muestra el menú de navegación (y en móvil un menú hamburguesa) y enlaces condicionados al estado de autenticación. \textbf{Por qué:} experiencia uniforme y una sola fuente de verdad para la navegación. La navegación entre pantallas se realiza con React Router y el hook \texttt{useNavigate}, evitando \texttt{window.location.href} para no provocar recargas completas y mantener el estado de la SPA.

\subsection{Flujos principales implementados}

\begin{itemize}
    \item \textbf{Importación de mazos desde MarvelCDB:} El usuario introduce una URL de mazo o selecciona uno desde la interfaz. El frontend obtiene los datos de la API pública de MarvelCDB, convierte las cartas al formato interno (resolviendo cada carta por \texttt{marvelcdb\_code} contra la base de datos propia para asignar aspecto y héroe correctos) y envía el mazo al backend con \texttt{POST /api/decks}. Se permite importar sin estar logueado: en ese caso \texttt{user\_id} queda a \texttt{NULL} en el backend, de modo que la funcionalidad esté disponible antes del registro.
    \item \textbf{Registro de partidas:} En \texttt{AddGamePage} el usuario elige un mazo existente o importa uno; en \texttt{ConfigureGamePage} selecciona villano, dificultad y resultado y envía la partida al backend. El backend persiste en \texttt{game\_configurations} y dispara el reentrenamiento del modelo en segundo plano, tal como se describió en el diseño del sistema de recomendación.
    \item \textbf{Recomendación IA:} La página de recomendación solicita al usuario villano y dificultad, llama al endpoint de recomendaciones del backend (que carga el modelo desde \texttt{MODEL\_PATH} si existe o aplica el criterio de respaldo) y muestra los mazos sugeridos. El frontend no entrena el modelo; solo consume la API y presenta los resultados.
    \item \textbf{Favoritos y comentarios:} Se implementaron las llamadas a los endpoints de favoritos y comentarios de mazos, con actualización de la UI (toasts, recarga de datos) tras cada acción. En el listado de favoritos se añadió una llamada a \texttt{getDeckById} cuando hace falta para mostrar correctamente el nombre del autor del mazo y evitar mostrar ``anónimo'' por falta de datos.
\end{itemize}

\subsection{Validación y feedback al usuario}

Los formularios de creación y edición de mazos validan límites de cartas (40--50), unicidad de nombres (normalizando mayúsculas/minúsculas) y campos obligatorios antes de enviar al backend. Las acciones destructivas (eliminar mazo, cerrar sesión) se confirman con diálogos. Los resultados de las operaciones se comunican mediante toasts (éxito, error, información). \textbf{Por qué:} reducir envíos inválidos y dar feedback inmediato sin depender solo de mensajes de error del servidor.

% -----------------------------------------------------------------------------
\section{Despliegue y entorno de producción}
% -----------------------------------------------------------------------------

El frontend se despliega en \textbf{Vercel} como aplicación estática generada por Vite; el backend se despliega en \textbf{Render} como servicio web. En esta sección se explica cómo y por qué se configuró la persistencia de la base de datos y del modelo en el servidor, que es crítico para que el sistema funcione correctamente en producción.

\subsection{Necesidad de persistencia en el servidor}

En Render, el sistema de archivos del contenedor es \textbf{efímero}: tras un reinicio o un nuevo despliegue, cualquier archivo creado en el directorio de trabajo se pierde. Si la base de datos SQLite y el archivo del modelo SVM (\,\texttt{.pkl}) se guardaran en ese sistema de archivos, se perderían los usuarios, mazos, partidas y el modelo entrenado en cada despliegue. \textbf{Por eso} se utiliza un \textbf{disco persistente} proporcionado por Render, montado en una ruta fija dentro del contenedor.

\subsection{Uso del directorio \texttt{/data} y variables de entorno}

En el backend desplegado en Render se configuró un \textbf{volumen persistente} montado en la ruta \texttt{/data}. Toda la información que debe sobrevivir a reinicios y redespliegues se almacena bajo ese directorio:

\begin{itemize}
    \item \textbf{Base de datos:} La variable de entorno \texttt{DB\_PATH} se define apuntando a un archivo dentro de \texttt{/data}, por ejemplo \texttt{/data/db.sqlite}. El backend (FastAPI) abre la conexión a SQLite usando esa ruta, de modo que las tablas \texttt{users}, \texttt{decks}, \texttt{game\_configurations}, \texttt{user\_favorites}, \texttt{deck\_comments} y \texttt{cards} se crean y persisten en ese archivo. En el arranque del servicio se ejecuta la inicialización de tablas (y, si corresponde, la importación inicial de cartas desde MarvelCDB) sobre la base apuntada por \texttt{DB\_PATH}, de forma que la primera vez se crea el esquema y en arranques posteriores se reutiliza la misma base.
    \item \textbf{Modelo de IA:} La variable \texttt{MODEL\_PATH} apunta a un archivo dentro de \texttt{/data}, por ejemplo \texttt{/data/saved\_models/villain\_svm\_model.pkl}. Cuando el backend reentrena el modelo (tras registrar partidas), guarda el archivo en esa ruta; cuando se recibe una petición de recomendación, carga el modelo desde \texttt{MODEL\_PATH} si existe. Así el modelo persiste entre reinicios y no es necesario reentrenar desde cero en cada despliegue.
\end{itemize}

\textbf{Por qué centralizar en \texttt{/data}:} Render permite asociar un único disco persistente por servicio y montarlo en una ruta elegida (\texttt{/data}). Usar siempre \texttt{DB\_PATH} y \texttt{MODEL\_PATH} en el código (y no rutas fijas en el directorio de trabajo) garantiza que en desarrollo se puede usar otra ruta (p. ej. local) y en producción se usa siempre el volumen persistente, sin hardcodear rutas del servidor.

\subsection{Frontend en producción}

El frontend se construye con \texttt{npm run build} (Vite) y los artefactos se sirven desde Vercel. Para que las peticiones en producción vayan al backend correcto, la URL base del API se configura mediante la variable de entorno \texttt{VITE\_API\_BASE\_URL} en el proyecto de Vercel, con el valor del servicio en Render (p. ej. \texttt{https://marveldcb-backend.onrender.com/api}). Así el código del frontend no contiene la URL de producción; solo usa la variable, lo que permite cambiar de entorno sin recompilar. La autenticación con Auth0 en producción requiere que en el panel de Auth0 estén registradas las URLs de callback y logout del dominio de Vercel, tal como se indicó en el capítulo de experimentos y pruebas.

\subsection{CORS y arranque en frío}

El backend incluye configuración CORS que permite orígenes del frontend desplegado (dominio de Vercel), de modo que el navegador permita las peticiones desde la SPA. Además, la lógica de arranque del backend comprueba que el directorio \texttt{/data} (o las rutas de \texttt{DB\_PATH} y \texttt{MODEL\_PATH}) existan antes de crear tablas o importar cartas, y maneja errores de archivo no encontrado para evitar fallos silenciosos cuando el volumen no está aún disponible en un arranque en frío.
